\chap Effects of casualties

One feature that separates \nns\ from other computational 
techniques is the fact that knowledge is distributed amongst the
weights in the network.
Hence, any network is redundant, and can sustain a certain amount of
injury without being overly affected.
Once again, just how much damage a \nn\ can take is something that
has not yet been widely investigated, but the simple experiments
detailed here, even though possibly not statistically sound, (only
one network was used,) provide an important insight into general
network behaviour.

The approach taken involved selection of a network, then damaging
the network and examining its behaviour for the same data set used
in the last chapter.
The network selected had the largest~$\varepsilon$ used,
$\varepsilon = 1.0$, and the smallest number of middle layer \neus,
$n_{\rm layer\ 1} = 15$.
This combination was chosen because it had the lowest error rate out
of any of the networks examined previously, and because it was
easiest to edit, with only a small number of middle layer \neus.

Casualties in the network took two forms.
First, random weights were damaged (set to zero) to
simulate the destruction of connections between \neus.
Various percentages of total connections were removed, namely~2.5\%,
5\%, 10\%, 20\%, 30\%, and finally~40\%.
The total error, (as before, $E$, the sum of the square of the
difference between actual and expected outputs,) and actual error
(the number of cases where the output layer \neus\ chose the wrong
pattern,) were calculated.

Second, individual \neus\ in the middle layer were ``damaged,'' that
is all the connections between their outputs and the top (output)
layer were set to zero.
The effect of the removal of between~1 and~7 out of~15 possible
\neus\ was examined, with total and actual error again being used.

We examine below these two types of casualties.
A summary of results is given at the end of this chapter.

\sec Casualties in Connections

FishNET, the \nn\ simulator package used in this thesis, stores
networks as ASCII files.
Hence, generating casualties in the network is a simple matter of
randomly picking a weight and setting it to zero.
This procedure was followed several hundred times, to produce
weight casualty rates of~0\%, 2.5\%, 5\%, 10\%, 20\%, 30\%, and~40\%.
The results of this damage appear in the table below.
Actual and total error are given for each percentage group.
$$
\vbox{\tabskip=0pt \offinterlineskip
\def\tablerule{\noalign{\hrule}}
\halign{\strut#& \vrule#\tabskip=1em plus2em&
 #\hfil& \vrule width1.5pt#&
 \hfil#\hfil& \vrule#&
 \hfil#\hfil& \vrule#&
 \hfil#\hfil& \vrule#&
 \hfil#\hfil& \vrule#&
 \hfil#\hfil& \vrule#&
 \hfil#\hfil& \vrule#&
 \hfil#\hfil& \vrule#\tabskip=0pt\cr
\tablerule
&&&&\multispan{13}\hfil Casualty Rate (weights)\hfil&\cr
\tablerule
&&{\ninerm Error}&&0\%&&2.5\%&&5\%&&10\%&&20\%&&30\%&&40\%&\cr
\noalign{\hrule height1.5pt}
&&{\ninerm Total}&&4.97&&5.11&&6.09&&6.60&&6.82&&7.97&&16.23&\cr
\tablerule
&&{\ninerm Actual}&&3&&3&&4&&4&&3&&4&&15&\cr
\tablerule
}}
$$
\centerline{Figure 15. Table of total and actual errors.}

These figures are graphed in figure 16.

\figure{16} {8} {Total and actual error versus \% weight casualties}

Some very interesting features emerge from this data.
With~2.5\% casualties, there is no appreciable difference in actual
or total error.
Once the number of casualties is doubled to~5\%, the total error
increases by~20\% and the actual error rate by~25\% (but as we shall see,
this is not an overly relevant figure).
However, for a further doubling of casualties to~10\%, there is only
an~8\% increase in total error, and doubling the damage
again causes a {\sl decrease\/} in actual error,
indicating that there is some random variation in actual error
results (so the previous increasing figure can be neglected).
From~0\% to~20\%, there is only a~37\% increase in total error,
and {\sl no\/} increase in actual error.
So it appears that network behaviour is not greatly modified at
casualty rates of up to~20\%.
By the time casualties reach~30\%, total and actual error are rapidly
on the increase.

Looking at the graph, we see that the error starts to increase
rapidly beyond casualty rates around~30\%.
Error rates of both kinds become exponentially worse, and the
network can be said to be no longer able to recognize patterns
meaningfully.
(Actual error has risen from~$3\over20$ or~$4\over20$ cases
at~$20\to30$\% casualties, to~$15\over20$ by~40\% casualties.)

From this data, we draw the conclusion that our $\varepsilon = 1.0$
and $n_{\rm layer\ 1} = 15$ network can sustain random injuries in up to
between~20\% and~30\% of weights, and still function satisfactorily.
This is quite incredible.

Larger networks (with larger intermediate layers) probably behave in
much the same way.
Of course, for the same {\sl number\/} of weight casualties the
effects will be less, but for the same {\sl percentage\/} of the
total weight population the behaviour will probably be the same.

\sec Casualties in Middle Layer Neurons

In this experiment, various numbers of middle layer \neus\ were
disabled, by setting the connecting weights between these \neus\ and
the output layer to zero.
This simulates the destruction of the individual processing
elements.
Here the effect of the removal of between~0 and~7 out of~15 middle
layer \neus\ is examined.

A table containing total and actual error data for a varying number
of \neus\ removed is shown in figure~17.
$$
\vbox{\tabskip=0pt \offinterlineskip
\def\tablerule{\noalign{\hrule}}
\halign{\strut#& \vrule#\tabskip=1em plus2em&
 #\hfil& \vrule width1.5pt#&
 \hfil#\hfil& \vrule#&
 \hfil#\hfil& \vrule#&
 \hfil#\hfil& \vrule#&
 \hfil#\hfil& \vrule#&
 \hfil#\hfil& \vrule#&
 \hfil#\hfil& \vrule#&
 \hfil#\hfil& \vrule#&
 \hfil#\hfil& \vrule#\tabskip=0pt\cr
\tablerule
&&&&\multispan{15}\hfil Number of Casualties (\neus)\hfil&\cr
\tablerule
&&{\ninerm Error}&&0&&1&&2&&3&&4&&5&&6&&7&\cr
\noalign{\hrule height1.5pt}
&&{\ninerm Total}&&4.97&&4.97&&5.27&&7.47&&8.39&&9.76&&10.44&&11.21&\cr
\tablerule
&&{\ninerm Actual}&&3&&3&&3&&6&&6&&9&&10&&7&\cr
\tablerule
}}
$$
\centerline{Figure 17. Table of total and actual errors.}

These figures are graphed in figure~18.

\figure {18} {5} {Total and actual errors versus number of layer 1
\neu\ casualties}

Another lot of interesting results comes from the table and graph.
With up to~2 \neus\ removed, there is no appreciable increase in
actual or total error.
With more than~2 removed, the actual error doubles and continues to
increase as more \neus\ are removed.
As the graph shows, the network rapidly decreases in usefulness as
the number of damaged \neus\ increases.

The data, then, shows us that the removal of a \neu\ has a more
dramatic effect as a straight percentage of \neus\ in the middle
layer than weight damage has.
If more than about~13\% $({2\over15})$ are damaged the network can't
really be used.
It is possible as well to look at the removal of \neus\ as a form of
weight damage.
By looking at it in this way, some remarkable conclusions can be
drawn about the representation of knowledge in the network, 
without even having to examine the distribution of weight
patterns in the network.

If one \neu\ out of~15 is damaged, then this is equivalent to the
systematic (instead of random) removal of~$1\over15$th of the total
\neus\ in the network.
Once the number of systematically removed weights
reaches more than~13\%, the performance degrades rapidly.
This fact reveals some important aspects of network behaviour.
Even though the percentage of weights damaged systematically is
lower, the effect is much greater than random weight damage.

Behaviour such as this indicates that although the knowledge of any
pattern is widely distributed among the weights in the network, the
selection of features is centralised into groups of \neus\ in the
middle layer.
Hence the injuring of intermediate layer \neus\ %
removes almost entirely, instead of
partially as for weights, a network's ability to recognise features
of a pattern, which in turn affects the output of pattern
recognition by the top layer.

For illustrative purposes, a mean time between failure analysis is
undertaken below to show just how reliable a \nn\ is with regard to failure
in intermediate layer \neus.

\sec Mean Time Between Failure Analysis

Angus [28] gives a simple calculation for mean time between 
failure (MTBF), which we use below for illustrative purposes.
From [28],
$$
{\displaystyle\theta_{\rm sys}} = 
	{\displaystyle{\sum_{i=0}^{n-k} {n\choose i} (\lambda/\mu)^i}\over
		\displaystyle{k\lambda {n\choose k} (\lambda/\mu)^{(n-k)}}},
$$
where the notation is explained as
$$
\halign{#\hfil&$\qquad$#\hfil\cr
$n$&number of nominally identical units\cr
$k$&minimum number of operating units for system to operate successfully\cr
$\theta_{\rm sys}$&mean (successful operating) time between system failures\cr
$\tau$&mean time to repair an individual unit\cr
$\theta$&mean time to failure for an individual unit\cr
$\lambda$&failure rate for an individual unit, $\lambda\equiv 1/\theta$\cr
$\mu$&repair rate for an individual unit, $\mu\equiv 1/\tau$\cr
}
$$
and
$$
{n\choose i} = { n! \over i!\,(n-i)!}.
$$

Imagine that a \nn\ is made up of individual, replaceable elements in
a hardware simulation, with a structure identical to the one 
used in this and the
previous chapter (3 layers, 195, 15, and~5 \neus).
So~2 out of the~15 \neus\ in the middle layer can be damaged and have
almost no effect on network behaviour.
Assume that we have some method of knowing that any
individual \neu\ (op-amp) has malfunctioned and can replace it
within~1 hour.
Also assume that the mean time to failure for each op-amp is
$10\,000$ hours of continuous operation.
Then, MTBF ($\theta_{\rm sys}$) for this network
is~$7.337004\times10^8$ hours, or~$83\,756$ years.
Quite a long time.
(These figures come from~$n=15$, $k=13$, $\lambda = {1\over10\,000}$,
and~$\mu={1\over1}$.)

\sec Summary

This chapter's results reveal some important facts about the
representation of knowledge in back-propagation \nn s, as well as
showing just how fault tolerant \nn s are.

Random removal of weights has shown that the network is quite immune
to failures of this kind.
In fact, between~20\% and~30\% of weights can be set to zero, with very
little effect upon the operation of the network.
This implies that the knowledge of all patterns learnt is spread
over all the weights in the network.

Removal of middle layer \neus\ (or systematic removal of weights)
reveals a more astounding fact.
It is widely recognised that \neus\ in the middle layer act as
feature extractors, reacting to particular parts of an input to
produce recognition of the appropriate pattern.
These results have shown that this same conclusion can be drawn by
removing intermediate layer \neus, and observing the increased
errors.

